\subsection{OLG Model 20: Search-and-Matching OLG}
We now solve an OLG model with search-and-matching in the labor market. From the household side this looks like an OLG model with exogenous labor, together with an infinite horizon firm problem, and a general equilibrium condition that determines the probability of finding a job (and filling a vacancy). From the perspective of VFI Toolkit, the household, firm, and general equilibrium conditions are all standard, but we need a 'custom model statistic' as one variables that goes into a general equilibrium condition is not a standard model aggregate. If you have never seen the simple Diamond-Mortensen-Pissarides search-and-matching model before, it is probably easier to look at that first before this example which combines the search-and-matching with an OLG model.

We start with the simplest search-and-matching setup. All non-employed households exert exogenous search effort. Firms create a 'vacancy' and if the match with a worker this vacancy is filled, creating a job. Workers and firms use Nash-Bargaining to determine the wage. To avoid interactions between capital and labor, it is assumed that the hiring of workers is done by intermediate firms, while production is done by a (representative) final goods firm. Thus, intermediate firms hire workers in the search-and-matching labor market, and then intermediate firms 'rent' labor to the final goods firm. The final goods firm has a standard Cobb-Douglas production function. Matches between non-employed households (who are all searching for a job) and firm vacancies are generated by a Cobb-Douglas matching function. As is standard in basic search-and-matching models, the household 'probability of finding a job' and the intermediate firm 'probability of filling a vacancy' are both proportional to the 'market tightness' which is defined as the ratio of the number of vacancies to the number of unemployed, where unemployed are the non-employed households that are searching for a job (which by assumption is here all the non-employed households, since search effort is exogenous).

Setting up search-and-matching in VFI Toolkit can be understood as involving two main aspects.

First, whether the household is employed or not is a markov state in the household problem, and the transition probability from non-employed to employed is given by the 'probability of finding a job'. Since the 'probability of finding a job' depends on the market tightness, and since the market tightness is determined in general eqm, it follows that the markov process in the household problem is determined in general eqm. To permit the markov process in the household problem to be determined in general eqm we must set it up as a function, which we do using \textit{vfoptions.ExogShockFn}.

Second, the search-and-matching adds two general eqm eqns. One of these is the markup of the wage over the marginal product of capital, which is determined by the Nash bargaining between workers and intermediate firms. This is a standard general eqm eqn so easy enough. The other general eqm eqn is the 'free entry condition' which says that the costs of (entering and) opening a vacancy must equal the benefits (otherwise more/less firms would enter to get the them equal). The challenge is that the benefits of entering depends on the value of hiring a worker, which depends on the value function of the firm and the distribution of the unemployed households (as these are the potential hires). We can set this up as a standard general eqm eqn, combined with a 'custom model statistic' for the expected value of the firm over the distribution of unemployed households. This is done using \textit{heteroagentoptions.CustomModelStats}.

To keep things simple, this model eliminates bequests and does not balance the funding for the pensions; so the model 'leaks' money due to accidental bequests that don't go to anyone, and has money 'injected' by people getting pensions for 'free'. We could add these as two more general eqm eqns, like in the other OLG models. I omit them here just so you can more easily see what search-and-matching involves in the code.

We assume that all households are either employed or non-employed, and that all non-employed households are exerting exogenous search effort. Because search-and-matching occurs randomly, and the probability of finding a job is determined in general eqm while the probability of job separation is exogenous, from the household perspective this looks like an exogenous markov process on employed/non-employed. As a result the household problem looks like a standard exogenous labor supply model, with an exogenous markov shock taking just two values, employed and non-employed. Households also have an idiosyncratic labor productivity, $z$.
\begin{eqnarray*}
 V(a,z,emp,j)&=& \max_{c,a'} \frac{c^{1-\sigma}}{1-\sigma} + \beta s_j E[V(a',z',emp',j+1)|z,emp] \\
&&\text{if working, }j<=Jr: c+a'= (1-\tau) w z emp + (1+r) a \\
&&\text{if retired, }j\geq Jr: c+a'=pension+ (1+r) a  \\
\end{eqnarray*}
Note that the interpretation of $emp$ is that $emp=1$ is employed and $emp=0$ is unemployed. We can solve this household problem in the usual manner. There is no decision variable, one endogenous state ($a$), and two exogenous markov states ($z$, $emp$).

The transition probabilities for the employment state are given by,
\begin{eqnarray*}
 Pr(emp'|emp) & emp'=0 & emp'=1 \\
 emp=0 & 1-p_{find} & p_{find} \\
 emp=1 & p_{sep} & 1-p_{sep}
\end{eqnarray*}
where $p_{sep}$ is an exogenous job seperation rate, and $p_{find}$ is a job finding probability that will be determined in general equilibrium.

There is a representative final goods firm with a Cobb-Douglas production function,
\begin{equation*}
 Y=K^{\alpha} L^{1-\alpha}
\end{equation*}
and we assume perfect competition to get the familiar conditions relating the rate of return on capital and the wage to the marginal products of capital and labor, respectively.
\begin{eqnarray*}
 h=(1-\alpha)A K^{\alpha}L^{-\alpha}
 r=\alpha A K^{\alpha-1} L^{1-\alpha} - \delta
\end{eqnarray*}
where $r$ is the return to capital, and $h$ is the rental price for labor services (which occurs in the competive market for labor services in which intermediate firms rent labor services to the final goods firm).

The labor market involves search-and-matching. Intermediate firms post vacancies, and we denote the number of vacancies as $V$. Non-employed households are all counted as unemployed (as they all exert exogenous search effort) and we denote the number of unemployed as $U$.\footnote{Techincally, $V$ is the mass of vacancies, and $U$ the mass of unemployed. Mass not number.} Matching occurs according to the Cobb-Douglas matching function,
\begin{equation*}
 M(U,V)=m U^{\gamma} V^{1-\gamma}
\end{equation*}
which leads to a job-finding probability of $p_{find}(\theta)=m \theta^{1-\gamma}$, and a job-filling probability of $q_{fill}=m \theta^{-\gamma}$, where $\theta = \frac{V}{U}$ is the market tightness.\footnote{You can easily derive this from the definitions of $p_{find}\equiv M(U,V)/U$ and $q_{fill}\equiv M(U,V)/V$. Is a direct result of the constant-returns-to-scale assumption implicit in the choice of Cobb-Douglas matching fn.}

Lastly we need to define the intermediate firm problem. Intermediate firms sit between household labor supply and final goods firm labor demand. Intermediate firms open vacancies, and then match in the labor market with households according to the job-filling probability, $q_{fill}$. Intermediate firms pay households a wage (per labor efficiency unit) of $w$, and rent this labor to the final goods firm for price $h$ (per labor efficiency unit), they make money on the difference between $h$ and $w$ (on the markup of $h$ over $w$, or to say the same thing the markdown of $w$ from $h$). Job-seperation occurs with exogenous probability $p_{sep}$.\footnote{My notation wants $q_{sep}$ to be the job-separation probility for the firm, and $p_{sep}$ to be the job-separation probability for the household, same as $p_{find}$ and $q_{fill}$ were used for job-finding and job-filling. But in equilibrium we know that $q_{sep}=p_{sep}$ (the number of jobs is the same viewed from both firm and household perspective, which is not true about job finding/filling as the vacancies and number of unemployed differ). So I just already impose the $q_{sep}=p_{sep}$ by using $p_{sep}$ directly in the intermediate firm problem.}, except that at the end of the period prior to retirement (in period $Jr-1$) the job separation rate is 1. If a worker dies, then this also results in job separation, and so the full job separation rate is $p_{sep}+sj-sj \, p_{sep}$ (assume that job seperation and death are independent shocks). So once a match occurs and a worker is hired, the value of this job to the intermediate firm is,
\begin{equation*}
 J(z,j)=(h-w) z + (1-p_{sep}-sj+sj p_{sep}) \frac{1}{1+r_{t+1}} E_z[J(z',j+1)|z]
\end{equation*}
where $J(\dot,Jr)=0$ captures the job-separation at retirement (implicitly, this is saying that $p_{sep}=1$ at age $j=Jr-1$).
notice that the per-period profit of the intermediate firm is $(h-w)z$, and then it discounts the future based on the interest rate and on the probability the job still exists next period. This leaves the issue of firms opening vacancies which lead to jobs. We assume that posting a vacancy costs $c$, and we already have that the probability of job-filling is $q_{fill}$, and that the value of the job is $J(z)$. Noting that the job occurs 'next period' and that the productivity of the worker being hired is drawn from the population of the non-employed (as these are the job seekers) we get that the free-entry condition is given by,
\begin{equation*}
 c=q_{fill} \frac{1}{1+r_{t+1}} \int_{z,j} E_z[J(z',j)|z] d\mu(emp=0,j)
\end{equation*}

The remaining issue is how the wage $w$ relates to the price of labor services $h$. We assume a Nash-inspired wage rule,
\begin{equation*}
 w=(1-\eta) \omega + \eta (h+ c \theta)
\end{equation*}
where $\eta$ is the worker's wage bargaining power.

To close the model we need to deal with the profits made by intermediate firms. We do this by distributing the profits as a lump-sum payment to households. The profit per-period made by the intermediate firms is $(h-w)z$, and so the intermediate firm sector has a profit of $\int_z (h-w)z d\mu(emp=1)$. This just requires one more simple general eqm eqn.

From the perspective of VFI Toolkit the key difference to previous models is the three additional general equilibrium conditions: (i) the differential between the wage received by households, and the wage paid by final-goods firms (which wage differential goes to the intermediate labor firms), (ii) the distribution of firm profits as a lump-sum to households, and (iii) the free-entry condition that the cost of opening a vacancy equals the expected value of the resulting job (which determines the equilibrium probability of finding a job, $p_{find}$, via market tightness).



The definition of a stationary equilibrium for this model is now,
\begin{definition}
 A stationary equilibrium is price $r$ (and $w$), and government policies $\tau$ and $pension$ such that
 \begin{enumerate}
  \item Given prices ($r$ and $w$) and government policies ($\tau$, $pension$), the household value function, $V$, and related policy function, $policy$, solve the household problem.
  \item The agents distribution evolves according to,
   \\  $\mu(a_{j+1},z_{t+1},e_{t+1},j+1)=\int \mathbb{I}_{a_{j+1}=policy_{a'}} \frac{\mu^{j+1}}{\mu^j} d\mu(a,z,e,j) \Gamma(z_{t+1}|z_t) Pr(\epsilon_{t+1}), \; j=1,2,...,J,\text{ and } \mu(:,:,:,1)=jequaloneDist$.
  \item The aggregates are based on household policies and agent distribution,
   \\ $L=\int \kappa_j exp(z+e) policy_h d\mu$
   \\ $K=\int a d\mu$
   \\ $PensionSpending=\int pension \mathbb{I}_{j\geq Jr} d\mu$
   \\ $AccidentalBeqLeft= \int policy_{a'}(1-s_j) d\hat{\mu}$
  \item Labor markets clear: $w=(1-\alpha)A K^{\alpha}L^{-\alpha}$
   \\ Capital markets clear: $r=\alpha A K^{\alpha-1} L^{1-\alpha} - \delta$
   \\ Pension system balance: $PensionSpending=\tau*w*L$
   \\ Accidental bequest balance: $AccidentBeq=AccidentalBeqLeft/(1+n)$
 \end{enumerate}
\end{definition}
\noindent note that the household problem is slightly different, and that the agent distribution (and hence all the integrals) is also now over the state space that includes the idiosyncratic exogenous shocks, namely over $(a,z,e,j)$.\footnote{When implementing we need to substitute for $Income$ with its expression in terms of the household variables (see the household problem), as otherwise the integral over the agent distribution does not make sense. I do not do so here purely because otherwise the equation gets a bit messy.} Note that the evolution of the agent distribution has been modified to also include the probabilities of the shocks.

As mentioned eariler, this model should be closed by cleaning up the accidental bequests (which currently leak out of the model) and funding the pensions (which are currently injected into the model). While the model is 'incomplete' without these they have been deliberately omitted here so that it is clearer in the codes how the search-and-matching is implemented.

Quick comment on related literature, essentially there are not many papers doing this. De la Croix, Pierrard \& Sneessens (2013) - Aging and pensions in general equilibrium: Labor market imperfections matter, does an OLG with search-matching, but they use one big representative household so not really relevant. Lugauer (2012) - Demographic Change and the Great Moderation in an Overlapping Generations Model with Matching Frictions, is one of the first examples. Cheron, Hairault \& Langot (2013) - Life-Cycle Equilibrium Unemployment, is the other. Very few other papers seem to do this, Graham \& Ozbilgin (2021) - Age, industry, and unemployment risk during a pandemic lockdown, does transition path in an OLG with search-and-matching, plus they have endogenous separation. What we are doing here is random search, Menzio, Telyukova \& Visschers (2016) - Directed search over the life cycle, does something similar but with directed search.

\subsection{OLG Model 21: Endogenous Search Effort}
We make one change to OLG Model 20, namely endogenizing search effort. This leads to one major conceptual change, namely the model now distinguishes employed, unemployed and inactive (previously, there was just employed and non-employed); with unemployed being those non-employed and actively searching (search effort greater than zero) and inactive being those non-employed but not searching (search effort zero). To keep the model tractable, it is assumed that the probability of finding a job is linear in the search effort (doubling search effort doubles the probability of finding a job).

In terms of what this means for setting up and solving the model, there are some minor changes to general equilibrium conditions and the definitions of parameters, but the major change is to the household problem. The household problem now includes the search effort as a decision and that this decision influences the transitions to employment which in turn means that the employed/non-employed variable is now a semi-exogenous state (not an exogenous markov state as before).

The household problem is now,
\begin{eqnarray*}
 V(a,semiz,emp,j)&=& \max_{d,c,a'} \frac{c^{1-\sigma}}{1-\sigma}-\psi(d) + \beta s_j E[V(a',semiz',emp',j+1)|z,emp] \\
&&\text{if working, }j<=Jr: c+a'= (1-\tau) w z emp + (1+r) a \\
&&\text{if retired, }j\geq Jr: c+a'=pension+ (1+r) a  \\
&& semiz'(d,semiz)
\end{eqnarray*}
where $d$ is assumed to take three possible values (0,1,2).

The transition probabilities for the employment semi-exogenous state are given by,
\begin{eqnarray*}
 Pr(emp'|emp) & emp'=0 & emp'=1 \\
 emp=0 & 1-d\,p_{find} & d\,p_{find} \\
 emp=1 & p_{sep} & 1-p_{sep}
\end{eqnarray*}
where $p_{sep}$ is an exogenous job seperation rate, and $d\, p_{find}$ is a job finding probability. $p_{find}$ is the probability of finding a job per unit of search effort, and will be determined in general equilibrium.

The final goods firm problem is unchanged. The only change the the matching function in the search-and-matching labor market is that $U$ is now defined as $\int d (emp=0) d\mu$ (previously it was $\int (emp=0) d\mu$). The problem of the intermediate firm is unchanged, except that the free-entry condition is now,
\begin{equation*}
 c=q_{fill} \frac{1}{1+r_{t+1}} \int_{z,j} E_z[J(z',j)|z] d\, d\mu(emp=0,j)
\end{equation*}
where we the integral with respect to the non-employed is now weighted by their search effort. The determination of wages is unchanged.




